% paper_12_algebraic_vee.tex — GeoVac Paper 12
% Algebraic Two-Electron Integrals on the Prolate Spheroidal Lattice
% Status: Draft, March 14, 2026
% Target: arXiv (physics.comp-ph / quant-ph)
\documentclass[aps,pra,twocolumn,superscriptaddress,longbibliography]{revtex4-2}

\usepackage{amsmath,amssymb,amsthm,graphicx,xcolor,bm,braket,dcolumn,booktabs}
\usepackage{hyperref}

\newtheorem{theorem}{Theorem}

% Numeric-registry citation form (CLAUDE.md Sec. 15).  Renders the literal
% only, so the PDF stays self-contained; the key is a foreign key into
% debug/qa/numeric_registry.py, which C21 recomputes and cross-checks.
\newcommand{\gvq}[2]{#2}

\begin{document}

\title{Algebraic Two-Electron Integrals on the\\
Prolate Spheroidal Lattice}
\thanks{Paper 12 in the Geometric Vacuum series}

\author{J.~Loutey}
\affiliation{Independent Researcher, Kent, Washington}

\date{March 14, 2026}

%% ========================================================================
%% ABSTRACT
%% ========================================================================

\begin{abstract}
Paper~11 established the prolate spheroidal lattice as the natural
geometry for diatomic molecules, solving H$_2^+$ to $0.0002\%$ energy
error (spectral Laguerre solver) with zero free parameters.  The extension to two electrons requires
electron-electron repulsion integrals $\langle\phi_i|1/r_{12}|\phi_j\rangle$
over the six-dimensional two-particle configuration space.  The
standard approach---numerical quadrature on a grid in prolate
spheroidal coordinates---converges slowly due to the Coulomb
singularity at $r_{12} = 0$ and saturates at $\sim$80\% of the
exact H$_2$ dissociation energy $D_e$ regardless of basis size.

We replace numerical integration with algebraic evaluation using
the Neumann expansion of $1/r_{12}$ in prolate spheroidal
coordinates, which decomposes each six-dimensional integral into a
finite sum of products of one-dimensional auxiliary integrals
computable by recurrence relations.  The resulting two-electron
integrals are exact within the Neumann truncation order $L$, with
no six-dimensional quadrature and no fitted parameters.  The one step
historically left to numerical integration---the log-singular
second-kind Legendre moment $B_l$---also has a closed form
(Sec.~\ref{sec:Bl_closed}): the only transcendental content is the
exponential integral $E_1$, the Euler--Mascheroni constant $\gamma$
and $\ln$, so \emph{in the $\sigma$ sector the reduction is fully
quadrature-free}.  (The $|m| \ge 1$ channels keep the same kernel;
generalising the same $A_l$, $B_l$, $X_l$ recurrences to associated
Legendre functions (Sec.~\ref{sec:azimuthal}) gives a moment-recurrence
engine that is \emph{stable} in the $\delta$ ($|m| = 2$) sector where
the differentiation route diverges, and---with the same closed-form
seeds---is likewise fully quadrature-free.)

A systematic convergence study on H$_2$ at $R = 1.4011$~bohr
demonstrates that the Neumann algebraic $V_{ee}$ yields 92.2\% of
the exact $D_e$ with 27 basis functions, compared to 79.6\% for
numerical $V_{ee}$ at the same basis size---a uniform 12--20
percentage point improvement across all basis sizes from 1 to 72
functions.  The Neumann result plateaus at 92.4\%, confirming that
$V_{ee}$ integration error has been eliminated.  \textbf{[MEASURED]}
The remaining 7.6\% gap is a basis-space limit, and we identify which
one.  The basis used here is $\varphi$-independent, so it spans only
$m_1 = m_2 = 0$ configurations; but a $^1\Sigma_g^+$ state constrains
the \emph{total} $M = m_1 + m_2$ to zero, not each $m_i$
separately, so the $\pi^2$, $\delta^2$, \ldots\ configurations---every
one of them $^1\Sigma_g^+$---are absent by construction.  Restoring
them in the same basis with the same Neumann kernel gives
$99.1\%$ of $D_e$ at $|m| \le 1$ (stability envelope
$99.0$--$99.1\%$;\ Sec.~\ref{sec:azimuthal}), a gain of $11.6$~mHa that
no growth along the $\sigma$ axis reproduces (near-tripling the
$\sigma$ basis is worth $0.34$~mHa).  \textbf{[MEASURED]} That $99.1\%$
is where the \emph{monomial} radial basis $\xi^j$ stalls---a Hankel
moment problem whose overlap matrix is past double precision, not a
structural ceiling.  Re-basing the same span onto an
orthogonal-polynomial family is an exact change of basis that removes
the linear dependence, and the energy then climbs monotonically and
variationally to $\gvq{p12_rebased_de_pct_aopt}{99.81}\%$ of $D_e$
($\gvq{p12_rebased_err_mha_aopt}{0.32}$~mHa, inside chemical accuracy) at a
larger truncation and the variationally optimal exponent
(Sec.~\ref{sec:recondition}).  \textbf{[WITHDRAWN]} This paper
previously attributed the gap to the electron-electron cusp requiring
non-analytic $r_{12}^{1/2}$ and $r_{12}\ln r_{12}$ terms, and read
that as motivation for hyperspherical coordinates.  Neither holds:
no such term appears in the calculation above, nor in the grid-based
prolate spheroidal treatment of Tao, McCurdy and Rescigno, which
reaches $0.05$~mHa in these same coordinates.
\end{abstract}

\maketitle

%% ========================================================================
%% I. INTRODUCTION
%% ========================================================================

\section{Introduction}
\label{sec:intro}

The Geometric Vacuum program identifies and discretizes the
\emph{natural geometry} of each quantum system---the coordinate
system where the Schr\"odinger equation separates---replacing
continuous integration with algebraic operations on quantum number
labels~\cite{loutey_paper0, loutey_paper1, loutey_paper7}.  For
atoms, the natural geometry is Fock's three-sphere
$S^3$~\cite{Fock1935}, where the Coulomb potential disappears into
the conformal metric and the Schr\"odinger equation becomes a pure
Laplace--Beltrami eigenvalue problem.  For diatomic molecules, the
natural geometry is the prolate spheroid, where the two-center
Coulomb problem separates into coupled ordinary differential
equations~\cite{loutey_paper11}.

Paper~11 implemented this principle for the one-electron
molecule H$_2^+$, achieving $0.0002\%$ energy error (spectral Laguerre
solver) with zero free parameters.  The extension to two electrons introduces a
qualitatively new challenge: the electron-electron repulsion
integral
\begin{equation}
  V_{ee}^{ij} = \left\langle\phi_i\left|\frac{1}{r_{12}}\right|\phi_j\right\rangle,
  \label{eq:vee_element}
\end{equation}
which couples the coordinates of both electrons and cannot be
evaluated by separation of variables.  In the prolate spheroidal
Hylleraas-type CI framework, these integrals are six-dimensional
(three coordinates per electron).  The $\sigma$-only basis adopted
below reduces the effective dimensionality to four;
Sec.~\ref{sec:gap} shows that this is a \emph{restriction} on the
configuration space, not a symmetry of the $^1\Sigma_g^+$ state, and
measures what it costs.

The standard approach computes Eq.~\eqref{eq:vee_element}
numerically, using the complete elliptic integral $K(k)$ to perform
the azimuthal average of $1/r_{12}$ and Gauss--Legendre quadrature
for the remaining four dimensions $(\xi_1, \eta_1, \xi_2, \eta_2)$.
This approach has two fundamental limitations:
\begin{enumerate}
  \item The Coulomb singularity at $r_{12} = 0$ (the
    electron-electron coalescence point) is poorly resolved by
    polynomial quadrature, requiring increasingly fine grids for
    convergence.
  \item The exponential decay $e^{-\alpha(\xi_1 + \xi_2)}$ of the
    basis functions extends to $\xi \to \infty$, requiring
    truncation of the semi-infinite radial domain.
\end{enumerate}
As we demonstrate in Section~\ref{sec:convergence}, these
limitations cause the numerical $V_{ee}$ to saturate at $\sim$80\%
of the exact H$_2$ dissociation energy, regardless of how many
basis functions are employed.

This paper replaces numerical quadrature with the Neumann
expansion of $1/r_{12}$ in prolate spheroidal
coordinates~\cite{Neumann1878, Roothaan1951, Ruedenberg1951},
which converts each $V_{ee}$ matrix element into a finite
sum of products of one-dimensional integrals.  These auxiliary
integrals---angular moments $C_l(p)$, exponential moments
$A_n(\alpha)$, and first- and second-kind Legendre moments
$A_l(p,\alpha)$ and $B_l(p,\alpha)$---are computed by recurrence
relations and closed forms.  The log-singular second-kind moment
$B_l$, once the sole step left to adaptive quadrature, is now itself
closed-form (Sec.~\ref{sec:Bl_closed}): a single log-moment against
$Q_0$ whose primitive reduces to $E_1$, $\gamma$ and $\ln$.  The
six-dimensional repulsion integral is thereby reduced, \emph{in the
$\sigma$ sector}, to a finite sum of algebraic operations on quantum
number labels with no numerical integration of any dimension;\
Sec.~\ref{sec:azimuthal} states what the $|m| \ge 1$ extension costs.

The structure of this paper is as follows.
Section~\ref{sec:neumann} derives the Neumann expansion and
its factorization for $\sigma$-state matrix elements.
Section~\ref{sec:auxiliary} defines the auxiliary integrals
and their recurrence relations.
Section~\ref{sec:implementation} describes the computational
pipeline.  Section~\ref{sec:convergence} presents the
convergence study comparing numerical and algebraic $V_{ee}$.
Section~\ref{sec:gap} diagnoses the origin of the remaining
7.6\% gap and measures its removal.  Section~\ref{sec:implications} discusses
implications for the GeoVac program.
Section~\ref{sec:conclusion} concludes.

%% ========================================================================
%% II. THE NEUMANN EXPANSION
%% ========================================================================

\section{The Neumann Expansion}
\label{sec:neumann}

\subsection{Expansion of \texorpdfstring{$1/r_{12}$}{1/r12} in
prolate spheroidal coordinates}

The Coulomb kernel $1/|\mathbf{r}_1 - \mathbf{r}_2|$ can be expanded
in prolate spheroidal coordinates $(\xi, \eta, \varphi)$ using the
Neumann expansion~\cite{Neumann1878, Morse1953}:
\begin{equation}
  \frac{1}{r_{12}} = \frac{2}{R}\sum_{l=0}^{\infty}\sum_{m=0}^{l}
    (2-\delta_{m0})\,(-1)^m\,(2l+1)
    \left[\frac{(l-m)!}{(l+m)!}\right]^{2}
    P_l^m(\xi_<)\,Q_l^m(\xi_>)\,
    P_l^m(\eta_1)\,P_l^m(\eta_2)\,
    \cos m(\varphi_1 - \varphi_2),
  \label{eq:neumann_full}
\end{equation}
where $\xi_< = \min(\xi_1, \xi_2)$, $\xi_> = \max(\xi_1, \xi_2)$,
$P_l^m$ and $Q_l^m$ are associated Legendre functions of the first
and second kind, and $R$ is the internuclear distance.
\textbf{[MEASURED]} Earlier versions of this paper printed
Eq.~\eqref{eq:neumann_full} without the $(-1)^m$ and $(2l+1)$
factors and with the factorial ratio unsquared.  That form does not
reproduce $1/r_{12}$---checked pointwise, it diverges---whereas the
expression above agrees with $1/|\mathbf{r}_1-\mathbf{r}_2|$ to
$2\times10^{-6}$ by $l = 10$ at well-separated $\xi$.  The error went
undetected because every calculation in \emph{earlier versions} of
this paper used only the $m = 0$ specialisation
Eq.~\eqref{eq:neumann_sigma}, which was derived independently---from
the Legendre completeness relation, as stated below---rather than
specialised from the printed general form.  Two of the three
discrepancies do vanish at $m = 0$;\ the $(2l+1)$ does not, and
Eq.~\eqref{eq:neumann_sigma} carries it because that independent
derivation supplies it.  Sec.~\ref{sec:azimuthal} is the first use of the
general form, which is why the error surfaced only now.

The basis of Eq.~\eqref{eq:basis_function} below carries no
$\varphi$ dependence, so for \emph{that basis} the azimuthal
integration $\int_0^{2\pi} d\varphi_1 \int_0^{2\pi} d\varphi_2$
retains only the $m = 0$ component:
\begin{equation}
  \left\langle\frac{1}{r_{12}}\right\rangle_\varphi
  = \frac{8\pi^2}{R}\sum_{l=0}^{\infty}(2l+1)\,
    P_l(\xi_<)\,Q_l(\xi_>)\,
    P_l(\eta_1)\,P_l(\eta_2),
  \label{eq:neumann_sigma}
\end{equation}
where $P_l \equiv P_l^0$ and $Q_l \equiv Q_l^0$ are ordinary
Legendre functions.  The factor $8\pi^2 = 2 \times 4\pi^2$ combines
the $2/R$ prefactor of the Neumann
expansion~\eqref{eq:neumann_full} with the $4\pi^2$ from
$\int_0^{2\pi}\!\int_0^{2\pi}\!\cos 0\,d\varphi_1\,d\varphi_2$.
The $(2l+1)$ factor arises from the Legendre polynomial
normalization in the angular completeness relation
$\int_{-1}^{1}P_l(\eta)\,P_{l'}(\eta)\,d\eta = 2\delta_{ll'}/(2l+1)$.

Equation~\eqref{eq:neumann_sigma} is the key identity for the
$\sigma$-only sector: it separates the two-electron Coulomb kernel
into a sum of products of functions of individual coordinates.
\textbf{[SCOPE]} It is \emph{not} the whole kernel for a
$^1\Sigma_g^+$ state.  Retaining only $m = 0$ is a restriction to
$\sigma$ configurations, not a consequence of $\Sigma$ symmetry: the
$\Sigma$ label fixes $M = m_1 + m_2 = 0$, which $\pi^2$ and
$\delta^2$ configurations satisfy with $m_2 = -m_1 \ne 0$.  The
terms with $m \ne 0$ in Eq.~\eqref{eq:neumann_full} are precisely
what couples those configurations to the $\sigma$ sector.
Sections~\ref{sec:gap} and~\ref{sec:azimuthal} measure the cost of
dropping them.  Each term in the sum
involves $P_l(\xi_<) Q_l(\xi_>)$, which couples the two radial
coordinates, and $P_l(\eta_1) P_l(\eta_2)$, which factorizes
completely.

\subsection{Matrix element factorization}

Consider two basis functions of the Hylleraas type in prolate
spheroidal coordinates:
\begin{equation}
  \phi_i(\xi_1, \eta_1, \xi_2, \eta_2)
  = \mathcal{N}_i\,e^{-\alpha(\xi_1 + \xi_2)}\,
    \xi_1^{j_i}\xi_2^{k_i}\,\eta_1^{l_i}\eta_2^{m_i}
    + (1 \leftrightarrow 2),
  \label{eq:basis_function}
\end{equation}
where the $(1 \leftrightarrow 2)$ exchange ensures bosonic
spatial symmetry for the singlet state, and
$(j_i, k_i, l_i, m_i)$ are non-negative integer quantum
numbers.  The exponential parameter $\alpha$ controls the
spatial extent; for all basis functions in a given calculation
we use a common $\alpha$ optimized variationally.

The $V_{ee}$ matrix element between $\phi_i$ and $\phi_j$ is
\begin{equation}
  V_{ee}^{ij} = \int\!\!\int
    \phi_i^*\,\frac{1}{r_{12}}\,\phi_j\;
    J(\xi_1,\eta_1)\,J(\xi_2,\eta_2)\,
    d\xi_1\,d\eta_1\,d\xi_2\,d\eta_2,
  \label{eq:vee_integral}
\end{equation}
where $J(\xi,\eta) = (R/2)^3(\xi^2 - \eta^2)$ is the prolate
spheroidal Jacobian (with the $2\pi$ azimuthal factors already
absorbed into Eq.~\eqref{eq:neumann_sigma}).

Substituting the Neumann expansion~\eqref{eq:neumann_sigma}
and expanding the product of two basis functions, each
unsymmetrized term contributes a product of the form
\begin{equation}
  e^{-2\alpha(\xi_1+\xi_2)}\,\xi_1^{P_1}\xi_2^{P_2}\,
  \eta_1^{Q_1}\eta_2^{Q_2},
\end{equation}
where $P_1 = j_i + j_j$, $P_2 = k_i + k_j$,
$Q_1 = l_i + l_j$, $Q_2 = m_i + m_j$ are the combined
powers from bra and ket.

The Jacobian factor $(\xi^2 - \eta^2)$ for each electron
expands the integral into four sub-terms with shifted powers:
\begin{align}
  &(+1):\; (P_1\!+\!2,\; P_2\!+\!2,\; Q_1,\;\; Q_2), \notag\\
  &(-1):\; (P_1\!+\!2,\; P_2,\;\;\;\; Q_1,\;\; Q_2\!+\!2), \notag\\
  &(-1):\; (P_1,\;\;\;\; P_2\!+\!2,\; Q_1\!+\!2,\; Q_2), \notag\\
  &(+1):\; (P_1,\;\;\;\; P_2,\;\;\;\; Q_1\!+\!2,\; Q_2\!+\!2).
  \label{eq:jacobian_expansion}
\end{align}

For each such sub-term with effective powers
$(\tilde{P}_1, \tilde{P}_2, \tilde{Q}_1, \tilde{Q}_2)$,
the matrix element factorizes as
\begin{equation}
  \frac{\pi^2 R^5}{8}\sum_{l=0}^{L}(2l+1)\,
  X_l(\tilde{P}_1, \tilde{P}_2, 2\alpha)\,
  C_l(\tilde{Q}_1)\,C_l(\tilde{Q}_2),
  \label{eq:factorized}
\end{equation}
where $C_l(q) = \int_{-1}^{+1}\eta^q\,P_l(\eta)\,d\eta$ are
angular moment integrals, and
\begin{equation}
  X_l(p_1, p_2, \alpha) = \int_1^\infty\!\!\int_1^\infty
    \xi_1^{p_1}\xi_2^{p_2}\,e^{-\alpha(\xi_1+\xi_2)}\,
    P_l(\xi_<)\,Q_l(\xi_>)\,d\xi_1\,d\xi_2
  \label{eq:Xl_def}
\end{equation}
is the two-dimensional radial integral coupling the two
electrons through the $P_l(\xi_<) Q_l(\xi_>)$ kernel.

The prefactor $\pi^2 R^5/8$ arises from
$(2/R) \times (R/2)^6 \times (2\pi)^2 = \pi^2 R^5/8$,
combining the Neumann prefactor $2/R$, the two Jacobians
$(R/2)^3$ each, and the azimuthal integration $4\pi^2$.

\subsection{Selection rules}

The angular integrals $C_l(q)$ obey a powerful selection rule:
\begin{equation}
  C_l(q) = 0 \quad\text{if } q < l \text{ or } (q - l)
  \text{ is odd}.
  \label{eq:selection_rule}
\end{equation}
This follows from the orthogonality of Legendre polynomials
and parity: $P_l(\eta)$ has parity $(-1)^l$, so
$\eta^q P_l(\eta)$ integrates to zero over $[-1,1]$ unless
$q + l$ is even, i.e., $(q - l)$ is even.  The condition
$q \geq l$ ensures that $\eta^q$ has enough powers to produce
a nonzero overlap with $P_l$.

For a basis with maximum $\eta$-powers $q_{\max}$, the
effective Neumann truncation is $l_{\mathrm{eff}} =
\min(L, q_{\max})$, where $L$ is the nominal truncation
order.  We use $L = 20$ throughout, which is more than
sufficient since our basis functions have $q_{\max} \leq 8$.

%% ========================================================================
%% III. AUXILIARY INTEGRALS
%% ========================================================================

\section{Auxiliary Integrals and Recurrence Relations}
\label{sec:auxiliary}

The factorized form~\eqref{eq:factorized} reduces each
$V_{ee}$ matrix element to products of three families of
one-dimensional integrals: $C_l(q)$, $A_n(\alpha)$ and
its Legendre generalizations, and $B_l(p, \alpha)$.  The
$C_l$ and $A_n$ families are computed by pure recurrence
relations; the log-singular second-kind moment $B_l$ has a
closed form (Sec.~\ref{sec:Bl_closed}), with adaptive quadrature
retained only as an independent numerical check.

\subsection{Angular moments \texorpdfstring{$C_l(q)$}{Cl(q)}}

\begin{equation}
  C_l(q) = \int_{-1}^{+1}\eta^q\,P_l(\eta)\,d\eta.
  \label{eq:Cl_def}
\end{equation}
The base cases are
\begin{equation}
  C_0(q) = \begin{cases} 2/(q+1) & q\text{ even}, \\ 0 & q\text{ odd},\end{cases}
  \qquad
  C_1(q) = \begin{cases} 0 & q\text{ even}, \\ 2/(q+2) & q\text{ odd}.\end{cases}
\end{equation}
Higher orders follow from the Legendre recurrence:
\begin{equation}
  (l+1)\,C_{l+1}(q) = (2l+1)\,C_l(q+1) - l\,C_{l-1}(q).
  \label{eq:Cl_recurrence}
\end{equation}
The recurrence $C_{l+1}(q)$ requires $C_l(q+1)$, so computing
$C_{l_{\max}}(q)$ requires workspace up to $q + l_{\max}$.

\subsection{Exponential moments \texorpdfstring{$A_n(\alpha)$}{An(alpha)}}

\begin{equation}
  A_n(\alpha) = \int_1^\infty\xi^n\,e^{-\alpha\xi}\,d\xi.
  \label{eq:An_def}
\end{equation}
The base case is $A_0(\alpha) = e^{-\alpha}/\alpha$, and the
upward recurrence is
\begin{equation}
  A_n(\alpha) = \frac{n\,A_{n-1}(\alpha) + e^{-\alpha}}{\alpha}.
  \label{eq:An_recurrence}
\end{equation}
For negative indices, $A_{-1}(\alpha) = E_1(\alpha)$
(the exponential integral) and the downward recurrence
\begin{equation}
  A_{-n-1}(\alpha) = \frac{e^{-\alpha} - \alpha\,A_{-n}(\alpha)}{n}
  \label{eq:An_negative}
\end{equation}
extends to arbitrary negative $n$.

\subsection{First-kind Legendre moments
\texorpdfstring{$A_l(p, \alpha)$}{Al(p, alpha)}}

\begin{equation}
  A_l(p, \alpha) = \int_1^\infty\xi^p\,e^{-\alpha\xi}\,
  P_l(\xi)\,d\xi.
  \label{eq:Al_def}
\end{equation}
The base cases are $A_0(p,\alpha) = A_p(\alpha)$ and
$A_1(p,\alpha) = A_{p+1}(\alpha)$, since $P_0(\xi) = 1$
and $P_1(\xi) = \xi$.  The Legendre recurrence gives
\begin{equation}
  (l+1)\,A_{l+1}(p,\alpha) = (2l+1)\,A_l(p+1,\alpha) - l\,A_{l-1}(p,\alpha).
  \label{eq:Al_recurrence}
\end{equation}

\subsection{Second-kind Legendre moments
\texorpdfstring{$B_l(p, \alpha)$}{Bl(p, alpha)}}
\label{sec:Bl_closed}

\begin{equation}
  B_l(p, \alpha) = \int_1^\infty\xi^p\,e^{-\alpha\xi}\,
  Q_l(\xi)\,d\xi.
  \label{eq:Bl_def}
\end{equation}
These integrals involve the Legendre function of the second
kind $Q_l(\xi)$, which has a logarithmic singularity at
$\xi = 1$:
\begin{equation}
  Q_0(\xi) = \tfrac{1}{2}\ln\frac{\xi+1}{\xi-1},
  \qquad Q_l(\xi) \sim -\tfrac{1}{2}\ln(\xi-1) + O(1)
  \quad\text{as }\xi\to 1^+.
  \label{eq:Ql_singularity}
\end{equation}
The log singularity is integrable (the integrand behaves as
$\ln(\xi-1)\,e^{-\alpha\xi} \to \ln(\xi-1)\,e^{-\alpha}$
near $\xi = 1$), but it defeats polynomial quadrature rules
such as Gauss--Laguerre, which assume smooth integrands.

Rather than integrate~\eqref{eq:Bl_def} numerically, we reduce it to
closed form.  Write $Q_l = P_l\,Q_0 - W_{l-1}$, where
$W_{l-1}(\xi) = \sum_{k=1}^{l} P_{k-1}(\xi)P_{l-k}(\xi)/k$ is a
polynomial and $Q_0(\xi) = \tfrac12\ln\frac{\xi+1}{\xi-1}$.  The
polynomial part contributes ordinary exponential moments $A_n(\alpha)$;
the entire transcendental content is carried by one family of
\emph{log-moments} against $Q_0$,
\begin{equation}
  B_l(p,\alpha) = \sum_n \big(\xi^p P_l\big)_n\,L_n(\alpha)
  \;-\; \sum_n \big(\xi^p W_{l-1}\big)_n\,A_n(\alpha),
  \qquad
  L_n(\alpha) = \int_1^\infty \xi^n\,Q_0(\xi)\,e^{-\alpha\xi}\,d\xi,
  \label{eq:Bl_closed}
\end{equation}
where $(\,\cdot\,)_n$ denotes the coefficient of $\xi^n$.  The
primitive $L_n$ is itself closed-form.  Splitting
$Q_0 = \tfrac12[\ln(\xi+1)-\ln(\xi-1)]$ and substituting $u=\xi\mp1$,
\begin{align}
  L_n(\alpha) &= \tfrac12\big(L^{+}_n - L^{-}_n\big), \notag\\
  L^{-}_n &= e^{-\alpha}\sum_{j=0}^{n}\binom{n}{j}\,
             \frac{j!\,(H_j-\gamma-\ln\alpha)}{\alpha^{\,j+1}},
  \label{eq:Lminus}\\
  L^{+}_0 &= \frac{e^{-\alpha}\ln 2}{\alpha} + \frac{e^{\alpha}}{\alpha}\,E_1(2\alpha),
  \qquad
  L^{+}_n = \frac{e^{-\alpha}\ln 2}{\alpha} + \frac{n}{\alpha}L^{+}_{n-1}
            + \frac{1}{\alpha}\Big[\textstyle\sum_{i=0}^{n-1}(-1)^i A_{n-1-i}(\alpha)
            + (-1)^n e^{\alpha}E_1(2\alpha)\Big],
  \label{eq:Lplus}
\end{align}
with $H_j$ the harmonic number, $\gamma$ the Euler--Mascheroni constant,
and $E_1$ the exponential integral.  The whole $B_l$ table then follows
from the $l = 0, 1$ seeds by the Legendre recurrence
\begin{equation}
  (l+1)\,B_{l+1}(p,\alpha) = (2l+1)\,B_l(p+1,\alpha) - l\,B_{l-1}(p,\alpha).
  \label{eq:Bl_recurrence}
\end{equation}
Thus the log-singular moment introduces \emph{no} numerical
integration: its minimal transcendental content is exactly
$\{E_1(2\alpha),\,\gamma,\,\ln\}$, the exponential-integral signature of
the Coulomb / $Q_0$ projection.  Adaptive
quadrature~\cite{QUADPACK} with subdivision near $\xi = 1$ is retained
only as an independent numerical check, agreeing with
Eqs.~\eqref{eq:Bl_closed}--\eqref{eq:Lplus} to
$\lesssim\!10^{-20}$ (backing:\
\texttt{tests/test\_paper12\_general\_m\_neumann.py}).

\subsection{Two-dimensional radial integral
\texorpdfstring{$X_l(p_1, p_2, \alpha)$}{Xl(p1, p2, alpha)}}

The integral~\eqref{eq:Xl_def} couples the two radial
coordinates through the $P_l(\xi_<) Q_l(\xi_>)$ kernel.
We split it into two ordered regions:
\begin{align}
  X_l(p_1, p_2, \alpha) &=
    \int_1^\infty\!\!\int_1^{\xi_2}
      \xi_1^{p_1}\xi_2^{p_2}\,e^{-\alpha(\xi_1+\xi_2)}
      P_l(\xi_1)\,Q_l(\xi_2)\,d\xi_1\,d\xi_2 \notag\\
    &\quad+ (1 \leftrightarrow 2).
  \label{eq:Xl_split}
\end{align}

The inner integral
$S_l(p,\alpha,\xi_0) = \int_1^{\xi_0}\xi^p\,e^{-\alpha\xi}\,
P_l(\xi)\,d\xi$
is the \emph{incomplete} first-kind Legendre moment.  Expanding
$\xi^p P_l(\xi) = \sum_j a_j\,\xi^j$ as a polynomial and
integrating each monomial by parts against $e^{-\alpha\xi}$:
\begin{equation}
  S_l(p,\alpha,\xi_0) = A_l(p,\alpha) - e^{-\alpha\xi_0}
  \sum_j a_j\sum_{k=0}^{j}\frac{j!}{(j-k)!\,\alpha^{k+1}}
  \,\xi_0^{j-k}
  + (\text{boundary at }\xi=1).
  \label{eq:ibp}
\end{equation}

Substituting into the outer integral and combining the
$e^{-\alpha\xi_0}$ factor from $S_l$ with the
$e^{-\alpha\xi_2}$ factor from the outer integrand yields
$e^{-2\alpha\xi_2}$ terms, which are evaluated as
$B_l(\cdot, 2\alpha)$:
\begin{equation}
  X_l = A_l(p_1)\,B_l(p_2) + A_l(p_2)\,B_l(p_1)
  - \text{(correction terms involving }B_l(\cdot, 2\alpha)\text{)}.
  \label{eq:Xl_formula}
\end{equation}
The full expression involves sums over the polynomial
coefficients of $\xi^{p_i} P_l(\xi)$ and is implemented
as described in Section~\ref{sec:implementation}.

%% ========================================================================
%% IV. IMPLEMENTATION
%% ========================================================================

\section{Implementation}
\label{sec:implementation}

\subsection{The algebraic \texorpdfstring{$V_{ee}$}{Vee} pipeline}

The computation of the $V_{ee}$ matrix proceeds in four stages:

\begin{enumerate}
  \item \textbf{Table precomputation.}  For a given exponential
    parameter $\alpha$ and basis quantum numbers, compute:
    \begin{itemize}
      \item The angular moment table $C_l(q)$ for $l \leq l_{\mathrm{eff}}$,
        $q \leq 2q_{\max} + 2$, via the recurrence~\eqref{eq:Cl_recurrence}.
        Cost: $O(l_{\max} \cdot q_{\max})$, negligible.
      \item The first-kind Legendre moments $A_l(p, 2\alpha)$ via
        the recurrence~\eqref{eq:Al_recurrence}.  Cost: $O(l_{\max} \cdot p_{\max})$.
      \item The second-kind Legendre moments $B_l(p, 2\alpha)$ and
        $B_l(p, 4\alpha)$ in closed form
        (Eqs.~\eqref{eq:Bl_closed}--\eqref{eq:Bl_recurrence}).
        Cost: $O(l_{\max} \cdot p_{\max})$ arithmetic, no quadrature;
        this replaced the former adaptive-quadrature seeding, which had
        dominated the $V_{ee}$ build ($\sim$40$\times$ faster, identical
        values).
    \end{itemize}
  \item \textbf{$X_l$ table.}  Compute
    $X_l(p_1, p_2, 2\alpha)$ for all
    $l \leq l_{\mathrm{eff}}$, $p_1, p_2 \leq p_{\max}+l_{\max}+2$,
    using Eq.~\eqref{eq:Xl_formula} with the precomputed
    $A_l$ and $B_l$ tables.  Cost: $O(l_{\max} \cdot p_{\max}^2
    \cdot l_{\max})$ arithmetic operations.
  \item \textbf{Matrix assembly.}  For each pair $(i, j)$ of
    basis functions, expand the bra$\times$ket product into
    unsymmetrized terms, apply the Jacobian
    expansion~\eqref{eq:jacobian_expansion}, and sum the
    Neumann series~\eqref{eq:factorized} using precomputed
    $C_l$ and $X_l$ tables.  Cost per element: $O(l_{\max})$
    (a single loop over $l$ with table lookups).
  \item \textbf{Symmetrization.}  Exploit $V_{ee}^{ij} = V_{ee}^{ji}$
    to compute only the upper triangle.
\end{enumerate}

The total cost scales as $O(N_{\mathrm{bf}}^2 \cdot l_{\max})$
for the matrix assembly, plus $O(l_{\max}^2 \cdot p_{\max}^2)$
for table precomputation.  For $N_{\mathrm{bf}} = 27$ and
$l_{\max} = 8$ (the effective truncation), the $V_{ee}$ matrix
is computed in $\sim$1~s.

\subsection{Hybrid Hamiltonian: Numba \texorpdfstring{$T + V_{ne}$}{T+Vne}
plus algebraic \texorpdfstring{$V_{ee}$}{Vee}}

The one-electron integrals (kinetic energy $T$ and nuclear
attraction $V_{ne}$) are computed analytically using integration
by parts for the kinetic energy derivatives and direct quadrature
for $V_{ne}$, accelerated by Numba JIT
compilation~\cite{Numba2015}.  A dedicated Numba kernel
\texttt{hamiltonian\_tvne\_p0\_kernel} computes $T + V_{ne}$
without $V_{ee}$, avoiding the expensive elliptic integral
computation that dominates the full numerical kernel.

The complete Hamiltonian is assembled as
\begin{equation}
  H = H_{T+V_{ne}}^{\text{(Numba)}} + V_{ee}^{\text{(Neumann)}},
  \label{eq:hybrid_H}
\end{equation}
where the first term uses grid-based quadrature (for the
non-singular $T$ and $V_{ne}$ integrals) and the second
uses the algebraic Neumann expansion (for the singular
$V_{ee}$ integrals).  This hybrid approach plays to the
strengths of each method: grids handle smooth integrands
efficiently, while the Neumann expansion handles the
Coulomb singularity exactly.

\subsection{Computational cost comparison}

For the $V_{ee}$ matrix alone:
\begin{itemize}
  \item \textbf{Numerical} (elliptic integral quadrature):
    $O(N_{\mathrm{bf}}^2 \cdot N_\xi^2 \cdot N_\eta^2)$.
    At $N_\xi = 30$, $N_\eta = 20$: $\sim$3~s for
    $N_{\mathrm{bf}} = 6$, scaling as $N_\xi^2 N_\eta^2$
    per element.
  \item \textbf{Algebraic} (Neumann expansion):
    $O(N_{\mathrm{bf}}^2 \cdot l_{\max} + l_{\max}^2
    \cdot p_{\max}^2)$, pure arithmetic throughout---the
    closed-form $B_l$ tables (Sec.~\ref{sec:Bl_closed}) removed the
    former quadrature entirely.
\end{itemize}

The algebraic approach is both faster (the closed-form $B_l$ made the
$V_{ee}$ build $\sim$40$\times$ faster than the earlier
quadrature-seeded path) and \emph{exact within the Neumann
truncation}, whereas the numerical approach always carries grid
truncation error.

%% ========================================================================
%% V. CONVERGENCE STUDY
%% ========================================================================

\section{Convergence Study}
\label{sec:convergence}

\subsection{Computational setup}

We compute the H$_2$ ground-state energy at $R = 1.4011$~bohr
(near-equilibrium) using a Hylleraas-type CI with basis
functions~\eqref{eq:basis_function}.  The basis is parameterized
by $(j_{\max}, l_{\max})$, which control the maximum powers
of $\xi$ and $\eta$ respectively.  The exponential parameter
$\alpha$ is optimized variationally at each basis size using
golden-section search over $\alpha \in [0.5, 2.5]$.

For each basis size, we compute the ground-state energy using
both numerical and Neumann $V_{ee}$:
\begin{itemize}
  \item \textbf{Numerical:} Full grid quadrature with elliptic
    integral $K(k)$ on a $30 \times 20$ grid ($N_\xi \times N_\eta$).
  \item \textbf{Neumann:} Algebraic $V_{ee}$ via
    Eq.~\eqref{eq:factorized} with $L = 20$.
\end{itemize}

The exact non-relativistic H$_2$ energy at this geometry is
$E_{\mathrm{exact}} = -1.174475$~Ha~\cite{Kolos1968}, giving
a dissociation energy $D_e = 0.174475$~Ha relative to two
separated hydrogen atoms ($E = -1.0$~Ha).

\subsection{Results}

Table~\ref{tab:convergence} presents the complete convergence
study.  The basis configurations range from a single function
($j = 0$, $l = 0$, $N = 1$) to 72 functions ($j = 3$, $l = 3$).

\begin{table*}[t]
\caption{Convergence of H$_2$ ground-state energy at
  $R = 1.4011$~bohr.  $N$ is the number of basis functions.
  $D_e\%$ is the fraction of the exact dissociation energy
  $D_e = 0.174475$~Ha recovered.  $\Delta E = E_{\mathrm{Neumann}}
  - E_{\mathrm{numerical}}$.
  \label{tab:convergence}}
\begin{ruledtabular}
\begin{tabular}{lrddddddd}
  \multicolumn{1}{c}{Config}
  & \multicolumn{1}{c}{$N$}
  & \multicolumn{1}{c}{$E_{\mathrm{num}}$ (Ha)}
  & \multicolumn{1}{c}{$D_e\%_{\mathrm{num}}$}
  & \multicolumn{1}{c}{$t_{\mathrm{num}}$ (s)}
  & \multicolumn{1}{c}{$E_{\mathrm{Neu}}$ (Ha)}
  & \multicolumn{1}{c}{$D_e\%_{\mathrm{Neu}}$}
  & \multicolumn{1}{c}{$t_{\mathrm{Neu}}$ (s)}
  & \multicolumn{1}{c}{$\Delta E$ (Ha)} \\
\hline
$j\!=\!0,\,l\!=\!0$ &  1 & -1.073192 & 41.9 &    1 & -1.107907 & 61.8 &    3 & -0.0347 \\
$j\!=\!1,\,l\!=\!0$ &  3 & -1.083438 & 47.8 &    8 & -1.115282 & 66.1 &    9 & -0.0318 \\
$j\!=\!1,\,l\!=\!1$ &  6 & -1.103485 & 59.3 &   32 & -1.129606 & 74.3 &   32 & -0.0261 \\
$j\!=\!2,\,l\!=\!1$ & 12 & -1.107960 & 61.9 &  152 & -1.132373 & 75.9 &  131 & -0.0244 \\
$j\!=\!2,\,l\!=\!2$ & 27 & -1.138836 & 79.6 &  767 & -1.160961 & 92.2 &  705 & -0.0221 \\
$j\!=\!3,\,l\!=\!2$ & 46 & -1.139329 & 79.8 & 2248 & -1.161162 & 92.4 & 1907 & -0.0218 \\
$j\!=\!3,\,l\!=\!3$ & 72 & -1.139792 & 80.1 & 6076 & -1.161304 & 92.4 & 5271 & -0.0215 \\
\end{tabular}
\end{ruledtabular}
\end{table*}

Three features of Table~\ref{tab:convergence} deserve emphasis:

\paragraph{1. Systematic improvement.}
The Neumann $V_{ee}$ yields a lower (better) energy than
numerical $V_{ee}$ at \emph{every} basis size, by 12--20
percentage points.  This is not a sporadic improvement from
fortuitous cancellation; it is a uniform advantage arising from
the elimination of grid truncation error in $V_{ee}$.

\paragraph{2. Numerical plateau at $\sim$80\%.}
The numerical $V_{ee}$ saturates at $\sim$80\% $D_e$ from
$N = 27$ onward: adding basis functions from 27 to 72 improves
the energy by only 0.5~percentage points (79.6\% $\to$ 80.1\%).
This saturation is caused by grid truncation error in the
$V_{ee}$ integrals, which introduces a systematic positive bias
in the electron-electron repulsion.  The basis is flexible
enough to represent a better wavefunction, but the inaccurate
$V_{ee}$ prevents it from being realized variationally.

\paragraph{3. Neumann plateau at 92.4\%.}
The Neumann $V_{ee}$ also plateaus, but at a much higher
level: 92.2\% at $N = 27$ and 92.4\% at $N = 72$.  Since the
Neumann $V_{ee}$ is exact (within the $L = 20$ truncation,
which is converged by the selection rule), this plateau is
\emph{not} caused by integration error.  It identifies a genuine
basis completeness limit, analyzed in Section~\ref{sec:gap}.

\subsection{The sweet spot: 27 basis functions}

The data reveal a clear efficiency optimum at $N = 27$
($j = 2$, $l = 2$): this configuration achieves 92.2\% $D_e$
in 705~s, whereas tripling the basis to $N = 72$ gains only
0.2~percentage points at 7.5$\times$ the computational cost.
The diminishing returns beyond $N = 27$ reflect the fact that
additional $\sigma$-type basis functions with higher $\xi$
and $\eta$ powers cannot access the missing physics---which is
not in the $\sigma$ sector at all
(Sections~\ref{sec:gap} and~\ref{sec:azimuthal}).

\subsection{Grid dependence of numerical \texorpdfstring{$V_{ee}$}{Vee}}

To confirm that the numerical plateau is grid-limited, we
performed a grid convergence study for the 6-function basis
($j = 1$, $l = 0$) at fixed $\alpha = 1.0$:

\begin{table}[h]
\caption{Grid convergence of numerical $V_{ee}$ for the
  6-function basis at $R = 1.4011$~bohr, $\alpha = 1.0$.
  \label{tab:grid_convergence}}
\begin{ruledtabular}
\begin{tabular}{cdd}
  Grid ($N_\xi \times N_\eta \times N_\phi$)
  & \multicolumn{1}{c}{$D_e\%$}
  & \multicolumn{1}{c}{$t$ (s)} \\
\hline
  $12\times 8\times 8$  & 66.0 &   1 \\
  $20\times 14\times 16$ & 76.2 &   6 \\
  $30\times 20\times 24$ & 81.4 &  45 \\
  $40\times 26\times 16$ & 83.7 &  83 \\
  $60\times 40\times 16$ & 85.7 & 447 \\
\end{tabular}
\end{ruledtabular}
\end{table}

With the same 6 basis functions, $D_e\%$ increases from 66\% to
85.7\% as the grid is refined from $12\times 8$ to $60\times 40$.
The convergence is slow ($\sim$1.3~percentage points per grid
doubling) and the cost grows quadratically.  In contrast, the
Neumann result for this basis is grid-independent at
61.8\% $D_e$ (the basis itself limits accuracy at 6 functions; the
$V_{ee}$ angular series terminates exactly, and the $B_l$ moment is
closed-form, so no quadrature enters).

%% ========================================================================
%% VI. DIAGNOSIS OF THE 7.6% GAP
%% ========================================================================

\section{Diagnosis of the Remaining 7.6\% Gap}
\label{sec:gap}

\subsection{What the plateau tells us}

The Neumann plateau at 92.4\% establishes that the $V_{ee}$
integration is converged: the Neumann series at $L = 20$ is
exact for our basis (the selection rule~\eqref{eq:selection_rule}
caps the effective $l$ well below 20).  The 7.6\% gap is
therefore a property of the \emph{basis itself}, not of the
integral evaluation.

The basis functions~\eqref{eq:basis_function} are products
of integer powers of $\xi$ and $\eta$ times a common
exponential $e^{-\alpha\xi}$.  These are smooth, analytic
functions of the single-electron coordinates, and---the point that
matters here---they carry no $\varphi$ dependence at all.

\textbf{[MEASURED]} The convergence study itself shows that the gap
is not reachable along the directions the basis does span.  Going
from $N = 27$ to $N = 46$ to $N = 72$ moves the energy
$92.2\% \to 92.4\% \to 92.4\%$:\ $\gvq{p12_sigma_growth_mha}{0.34}$~mHa for $2.7\times$ the
functions.  Whatever is missing, more $\sigma$ functions do not
supply it.  Sections~\ref{sec:azimuthal_diagnosis}
and~\ref{sec:azimuthal} identify it and remove it.

\subsection{The electron-electron cusp is real, but it is not this
gap}
\label{sec:azimuthal_diagnosis}

The exact two-electron wavefunction $\Psi(\mathbf{r}_1,
\mathbf{r}_2)$ satisfies the Kato cusp condition~\cite{Kato1957}
at the electron-electron coalescence point $r_{12} = 0$:
\begin{equation}
  \left.\frac{\partial\bar{\Psi}}{\partial r_{12}}\right|_{r_{12}=0}
  = \frac{1}{2}\,\bar{\Psi}(r_{12}=0),
  \label{eq:kato}
\end{equation}
where $\bar{\Psi}$ is the spherically averaged wavefunction.
This cusp produces a discontinuity in the first derivative of
$\Psi$ with respect to $r_{12}$ and, more fundamentally,
introduces non-analytic dependence on $r_{12}$.

The Fock expansion~\cite{Fock1954, Morgan1986} of the exact
wavefunction near the two-electron coalescence point has the form
\begin{equation}
  \Psi = \sum_{n=0}^\infty\sum_{k=0}^{\lfloor n/2\rfloor}
    R^{n/2}\,(\ln R)^k\,f_{nk}(\hat{\Omega}),
  \label{eq:fock_expansion}
\end{equation}
where $R = (r_1^2 + r_2^2)^{1/2}$ is the hyperradius and
$\hat{\Omega}$ represents the hyperangular coordinates.  The
$R^{1/2}$ and $R\ln R$ terms are fundamentally non-analytic:
no finite sum of polynomial$\times$exponential functions in
prolate spheroidal coordinates can represent them \emph{exactly}.

\textbf{[WITHDRAWN]} Earlier versions of this paper concluded from
this that the cusp is what caps the calculation at 92.4\%.  It is
not, and the argument was never sound:\ non-analyticity implies slow
convergence of a partial-wave expansion, not a floor at 7.6\%.  Two
calculations settle it.  Tao, McCurdy and
Rescigno~\cite{tao_mccurdy_rescigno2010}, working in these same
prolate spheroidal coordinates with a polynomial angular basis and no
$r_{12}$ dependence anywhere in it, reach
$-1.17442$~Ha---$0.05$~mHa from exact---at angular truncation
$l_{\max} = 6$.  (Their odd-$m$ radial functions do carry an explicit
$(\xi^2-1)^{1/2}$, which they introduce for the non-analytic behaviour
at $\xi = 1$;\ that is a one-electron boundary condition on the focal
axis, the analogue of the $\sin\theta$ behaviour of ordinary spherical
harmonics at odd $m$, and is unrelated to $r_{12}$.  The point stands
as stated:\ nothing in their basis represents the \emph{interelectronic}
non-analyticity.)  And the calculation reported in
Sec.~\ref{sec:azimuthal} below, run in the basis of
Eq.~\eqref{eq:basis_function} with the same Neumann kernel and
likewise no $r_{12}$ dependence, reaches $99.1\%$.  Neither could
exist if representing the cusp were the binding constraint at this
accuracy.  What the cusp does cost is the slow tail of the
partial-wave series;\ at $|m| \le 1$ and $l \le 3$ that residue is
under one percent of $D_e$.

\subsection{Why \texorpdfstring{$r_{12}$}{r12} basis functions
help only partially}

The traditional approach to the cusp problem is to include
explicit $r_{12}$ dependence in the basis, as in the classic
James--Coolidge~\cite{James1933} and Kolos--Wolniewicz~\cite{Kolos1968}
wavefunctions.  We tested this by extending the basis with
$r_{12}^p$ ($p = 1, 2$) factors.  The results are instructive:
with 18 functions including $r_{12}^2$ and angular terms
($l \leq 1$), the numerical $V_{ee}$ reaches 86.8\% $D_e$
at the $20\times 14\times 16$ grid---but this \emph{also}
saturates with grid refinement.  The $r_{12}$ factors improve
the wavefunction's ability to represent the cusp, but the
$V_{ee}$ integrals involving $r_{12}$ are even harder to
evaluate numerically (they become seven-dimensional integrals
with an additional radial singularity).

The Neumann algebraic approach at 92.2\% with $p = 0$ basis
functions outperforms the $r_{12}$ + numerical approach at
86.8\% with $p \leq 2$ basis functions.  What that comparison
establishes is narrower than it looks:\ \emph{an exact $V_{ee}$ beats
$r_{12}$ basis functions carrying grid error}.  It says nothing about
the cusp in either direction.  (Earlier versions drew the stronger
conclusion that ``the cusp does not need $r_{12}^p$ basis
functions---it needs exact $V_{ee}$'';\ that goes with the rest of the
withdrawn cusp reading, and
Sec.~\ref{sec:azimuthal_diagnosis} gives the cusp's actual cost here.)  The $r_{12}$ basis functions are
a workaround for poor $V_{ee}$ integration, not a solution
to the underlying representation problem.

\subsection{The missing channels}

\textbf{[MEASURED]} The 7.6\% gap is a \emph{configuration-space}
limitation, and an elementary one.  A $^1\Sigma_g^+$ state requires
the \emph{total} azimuthal quantum number $M = m_1 + m_2$ to vanish.
It does not require $m_1 = m_2 = 0$.  The configurations
$\pi^2$ ($m_1 = +1$, $m_2 = -1$), $\delta^2$ ($m_1 = +2$,
$m_2 = -2$), and so on all have $M = 0$ and are all
$^1\Sigma_g^+$;\ they carry the \emph{angular} part of the electron
correlation, the left--right and in--out parts being what a
$\sigma$-only space already describes.  The basis
Eq.~\eqref{eq:basis_function} is $\varphi$-independent and therefore
excludes every one of them, and Eq.~\eqref{eq:neumann_sigma} drops
exactly the kernel terms that would couple them in.  The two
restrictions are consistent with each other, which is why the
calculation converges cleanly to a wrong answer.

Section~\ref{sec:azimuthal} restores the channels and measures the
result:\ $92.4\% \to 99.1\%$ of $D_e$.  What remains after that is
ordinary basis incompleteness---higher $l$, higher $n$, and the
$\delta$ channels.  That remainder is ${\approx}1.6$~mHa in total, of
which the $\delta$ channels account for ${\approx}0.5$~mHa on an
independent Gaussian-basis estimate;\ the balance is not itemised
here.

For completeness, the interelectron coordinate in prolate
spheroidals is a function of all five relative variables:
\begin{equation}
  r_{12}^2 = \frac{R^2}{4}\bigl[
    (\xi_1^2-1)(1-\eta_1^2) + (\xi_2^2-1)(1-\eta_2^2)
    - 2\sqrt{(\xi_1^2-1)(1-\eta_1^2)(\xi_2^2-1)(1-\eta_2^2)}
      \,\cos(\varphi_1-\varphi_2)
    + (\xi_1\eta_1 - \xi_2\eta_2)^2
  \bigr].
  \label{eq:r12_prolate}
\end{equation}
(The $\cos(\varphi_1-\varphi_2)$ factor was omitted in earlier
versions of this equation, which therefore held only at
$\varphi_1 = \varphi_2$.)

Note the $\cos(\varphi_1-\varphi_2)$ dependence:\ the
interelectron distance is not a function of
$\xi_1, \eta_1, \xi_2, \eta_2$ alone.  A $\varphi$-independent basis
cannot resolve it, which is the same restriction stated in
coordinates.

\textbf{[WITHDRAWN]} Earlier versions of this paper closed this
section by identifying hyperspherical coordinates
$(\rho, \theta, \hat{\Omega})$ as the natural geometry the gap
called for, on the grounds that the coalescence cusp has a simple
representation there.  That inference rested on the cusp diagnosis
withdrawn above and does not survive it.  Hyperspherical coordinates
remain the right setting for genuine three-body coalescence
problems, and Paper~13~\cite{loutey_paper13} uses them for helium
on independent grounds;\ what is withdrawn is the claim that
\emph{this} calculation's residual motivates them.

%% ========================================================================
%% VI-B. RESTORING THE AZIMUTHAL CHANNELS
%% ========================================================================

\section{Restoring the Azimuthal Channels}
\label{sec:azimuthal}

\textbf{[MEASURED]} We extend the basis of
Eq.~\eqref{eq:basis_function} with an azimuthal quantum number,

\begin{equation}
  u_{j l \mu}(\xi,\eta) =
    \xi^{\,j}\,\eta^{\,l}\,
    (\xi^2-1)^{\mu/2}(1-\eta^2)^{\mu/2}\,e^{-\alpha\xi},
  \qquad
  \Phi_\mu = u\,u'\,\cos\mu(\varphi_1-\varphi_2),
  \label{eq:basis_mu}
\end{equation}
pairing $m_1 = +\mu$ with $m_2 = -\mu$ through
$\cos\mu(\varphi_1-\varphi_2)$---the $M = 0$, $^1\Sigma_g^+$
combination, which for $\mu = 1$ is the familiar
$\pi_x\pi_x + \pi_y\pi_y$.  At $\mu = 0$ the factor $u$ reduces to the one-electron factor of
Eq.~\eqref{eq:basis_function} and $\Phi_0$ to its symmetrised
product, so the $\mu = 0$ sector is exactly the original basis.  The overlap, kinetic and
nuclear-attraction matrices remain \emph{exact}---every integrand is
a polynomial in $\xi$ and $\eta$ times $e^{-2\alpha\xi}$, evaluated
against the moments $A_n$ of Sec.~\ref{sec:auxiliary} and
elementary $\eta$ moments---and all three are diagonal in $\mu$.
$V_{ee}$ uses the full kernel Eq.~\eqref{eq:neumann_full};\ its
$m \neq 0$ terms are what couple different $\mu$.

The $\eta$ integrals remain exact polynomial moments, because the
kernel's $(1-\eta^2)^{m/2}$ and the basis's $(1-\eta^2)^{\mu/2}$
always combine to an integer power.  The sum over $l$ terminates
exactly:\ integrating by parts $m$ times (the boundary terms vanish
since $(1-\eta^2)^{s}$ has a zero of order $s \ge m$ at
$\eta = \pm1$) shows the moment is identically zero for
$l > Q + 2s - m$, with $Q$ the combined $\eta$ power and
$s = (\mu_i + \mu_j + m)/2$.  That rule must be imposed
explicitly rather than left to cancellation:\ the Legendre
derivative coefficients reach $\sim\!10^{10}$, and their
floating-point residue multiplied by the radial integral produced a
$-3\times10^{8}$~Ha artifact before the cutoff was enforced.

\begin{table}[t]
\caption{H$_2$ at $R = 1.4011$~bohr with the azimuthal channels
  restored.  Same basis family, same Neumann kernel, $\alpha$
  scanned, canonical orthogonalization throughout --- eigenvectors of
  $S$ with eigenvalue below $10^{-11}\lambda_{\max}$ are discarded, and
  the surviving dimension is given in parentheses where it differs
  from $N$.  $D_e^{\rm exact} = 0.174475$~Ha.}
\label{tab:azimuthal}
\begin{ruledtabular}
\begin{tabular}{crdrd}
  \multicolumn{1}{c}{$(j_{\max}, l_{\max})$}
  & \multicolumn{1}{c}{$N_{\sigma}$}
  & \multicolumn{1}{c}{$D_e\%$ ($\mu = 0$)}
  & \multicolumn{1}{c}{$N_{|m|\le1}$}
  & \multicolumn{1}{c}{$D_e\%$ ($|m| \le 1$)} \\
\hline
$(1,1)$ &   6 & 74.19 &  12 & 81.63 \\
$(2,1)$ &  12 & 75.87 &  24 & 83.47 \\
$(2,2)$ &  27 & 92.25 &  54 & 98.96 \\
$(3,2)$ &  46 & 92.37 &  92 (76) & 99.00 \\
$(3,3)$ &  72 (65) & \gvq{p12_sigma_only_de_pct}{92.42} & 144 (115) & \gvq{p12_azimuthal_de_pct}{99.09} \\
\end{tabular}
\end{ruledtabular}
\end{table}

Table~\ref{tab:azimuthal} gives the result.  At the largest basis
the azimuthal channels are worth $+11.64$~mHa, taking
$\gvq{p12_sigma_only_de_pct}{92.42}\% \to \gvq{p12_azimuthal_de_pct}{99.1}\%$ of $D_e$.  The $\mu = 0$ column reproduces the
Neumann column of Table~\ref{tab:convergence} to
$161$, $1.6$, $1.0$, $6.7$ and $58$~$\mu$Ha.

Three controls make the reading unambiguous.  \emph{First}, the
gain is not a basis-count effect:\ tripling the basis along the
$\sigma$ axis is worth $\gvq{p12_sigma_growth_mha}{0.34}$~mHa
(Sec.~\ref{sec:gap}), while
opening the azimuthal axis at fixed $(j_{\max}, l_{\max})$ is worth
$\gvq{p12_azimuthal_gain_mha}{11.6}$~mHa.  \emph{Second}, an independent
route agrees.  A full CI
in a Cartesian Gaussian basis---different functions, different
integrals, different code---restricted to $m = 0$ orbitals converges
to $92.34\%$, within $0.2$~mHa of this paper's $\sigma$-only value
in a completely unrelated basis;\ releasing $|m| \le 1$ takes it to
$99.10\%$, against $99.1\%$ here---agreement well inside the
stability envelope quoted below, and not to be read as a
two-decimal coincidence.  \emph{Third}, the general-$m$
kernel of Eq.~\eqref{eq:neumann_full} was checked pointwise against
$1/|\mathbf{r}_1-\mathbf{r}_2|$, agreeing to better than $10^{-6}$ at
well-separated $\xi$ by $l = 10$.  That figure is a point evaluation, not
a uniform bound:\ the Neumann series converges slowly when the two $\xi$
are close, and the second-kind recursion used here degrades at larger
$l$, so nearby geometries and higher truncations are worse.  What the
check establishes is the \emph{form} of the prefactor---the superseded
one errs by nine orders of magnitude at the same point---not an accuracy
claim for the kernel.

Backing:\ \texttt{tests/test\_paper12\_azimuthal\_channels.py},
whose guards are fire-tested against the specific wrong answers they
exclude (a killed azimuthal coupling, the superseded kernel prefactor,
a disabled canonical orthogonalization, and a truncation cap that would
hide the selection rule).

\textbf{[MEASURED]} Both limits stated in earlier versions of this
section have since been removed.  The $\mu > 0$ radial integrals are
now available from the $A_l$, $B_l$, $X_l$ recurrences of
Sec.~\ref{sec:auxiliary} generalised to the associated Legendre
functions (module \texttt{geovac/neumann\_vee\_general\_m.py}).  The
first-kind moment $A_l^{m,s}(p) = \int_1^\infty \xi^p
(\xi^2-1)^s e^{-c\xi}\,[\,\mathrm{d}^m P_l/\mathrm{d}\xi^m\,]\,
\mathrm{d}\xi$ is a closed-form polynomial contraction against the
$A_n$; the second-kind moment $B_l^{m,s}$---whose integrand is
singular as $(\xi-1)^{-m}$ at $\xi = 1$, regularised by the weight
$(\xi^2-1)^s$ kept \emph{intact} rather than expanded into
$\xi$-powers (the bare $\xi^p\,\mathrm{d}^m Q_l$ moment diverges for
$m \ge 1$)---is built by the associated-Legendre forward $l$-recurrence
$(l-m+1)B_{l+1}(p) = (2l+1)B_l(p+1) - (l+m)B_{l-1}(p)$, seeded at
$l = m, m+1$; the ordered-$\xi$ integral $X_l^{m,s}$ then follows from
the same integration-by-parts split as the $\sigma$ case.  The engine
reduces to \texttt{geovac.neumann\_vee} at $m = 0$ to $\sim\!10^{-9}$
elementwise, matches a high-precision reference for $X_l^{m,s}$ to
better than $10^{-9}$ through $l = 8$, $m \le 4$, and reproduces the
grid result at $(2,2)$, $|m| \le 1$ to the last printed digit.  Its
decisive property is \emph{stability} in the $\delta$ ($|m| = 2$)
sector:\ the fourth-derivative $Q_l$ cancellation near $\xi = 1$ drives
the differentiation route to a non-variational $\approx -14.5$~Ha,
whereas the recurrence build is variational and supplies the $\delta$
channels from the native prolate route---a $+0.25$~mHa gain at
$(2,2)$, consistent with the $\approx 0.5$~mHa Gaussian estimate (the
$(2,2)$ $\delta$ subspace being the smaller).  The engine is
\emph{fully quadrature-free}:\ the low-$l$ $B_l^{m,s}$ seeds are
closed-form by the same construction as the $\sigma$ sector
(Sec.~\ref{sec:Bl_closed}).  The obstruction that had kept these seeds
numerical---the individual pole terms $d^{k}Q_0 \sim (\xi-1)^{-k}$ of
$d^m Q_l$ have divergent moments---dissolves because $s \ge m$ always
holds ($s = (\mu_i+\mu_j+m)/2$ with $m \in \{\mu_i+\mu_j,
|\mu_i-\mu_j|\}$): the intact weight
$(\xi^2-1)^s = (\xi-1)^s(\xi+1)^s$ then \emph{fully polynomialises}
every pole ($k = m-a \le m \le s$), leaving a single log-moment against
$Q_0$ whose primitive is Eqs.~\eqref{eq:Bl_closed}--\eqref{eq:Lplus}.
The moment tables are still built in extended precision (the associated
Legendre coefficients cancel over several orders), but no numerical
integration remains.  Backing:\
\texttt{tests/test\_paper12\_general\_m\_neumann.py}.

\textbf{[MEASURED]} Finally, a caution about the $\sigma$ column
itself.  The overlap matrix of this basis family is strongly
linearly dependent at high powers and a single $\alpha$:\ measured
$\mathrm{cond}(S) = 2.6\times10^{14}$ at $(j,l) = (3,3)$, $\mu = 0$
(the $N = 72$ basis of Table~\ref{tab:convergence}), rising to
$2.0\times10^{16}$ once $\mu = 1$ doubles the basis.  Beyond
$\sim\!10^{16}$ a direct generalised eigensolve is not
trustworthy---at $N = 144$ it returned $-79$~Ha---so
Table~\ref{tab:azimuthal} uses canonical orthogonalization
throughout.

\textbf{[MEASURED]} Canonical orthogonalization \emph{bounds} that
pathology;\ it does not remove it, and we state the envelope rather
than leave the reader to assume a converged four-figure number.  At
the $|m| \le 1$, $N = 144$ basis the solver still returns
non-variational values at a \emph{majority} of $\alpha$.  Measured at
the stated threshold over $\alpha \in [0.90, 1.30]$, six of nine grid
points fail.  The failures are not marginal:\ $\alpha = 1.15$ and $1.20$
return $-4.1$ and $-8.1$~Ha, $\alpha = 0.95$ returns $-277$~Ha, and
$\alpha = 1.00$---the natural default, and the value a reader is most
likely to try first---returns $-64.3$~Ha.  The three that survive give $99.03$, $99.09$ and $99.05\%$ at
$\alpha = 1.10$, $1.25$, $1.30$.  So the reported value is the best
\emph{variational} point of an $\alpha$ scan, and points failing the
variational test are discarded---a selection the reader should see, not
infer.  The value also moves with the discard threshold:\ at
$\alpha = 1.25$ it reads $99.14\%$ at $10^{-12}$, $99.09\%$ at
$10^{-11}$, $98.99\%$ at $10^{-10}$ and $98.41\%$ at $10^{-8}$,
and is non-variational below $\sim\!10^{-12}$.  The honest
statement is therefore $99.0$--$99.1\%$, with
$\gvq{p12_azimuthal_de_pct}{99.09}\%$ the best
variational value at the stated parameters.

\textbf{[SCOPE 2026-09-14, corrected twice]}
Two earlier versions of this passage overstated how stable that number is.
The first called the conclusion ``insensitive to the choice''
[retracted 2026-09-14: p12-envelope-insensitive] and then quoted the figure
that refutes it.  The second replaced it with a range, ``$95.5$--$99.1\%$
across the grid'', and called the qualitative conclusion ``robust across the
grid'' [retracted 2026-09-14: p12-grid-floor-955].  Both are wrong, and the
second is wrong in the direction that flatters the result.

$95.5\%$ is the cell at $\alpha = 1.10$, threshold $10^{-8}$.  It is the
weakest variational value only if $\alpha$ is restricted to $\ge 1.10$---which
excludes precisely the low-$\alpha$ half this paragraph has just identified as
pathological.  Over the full declared range $\alpha \in [0.90, 1.30]$ the
weakest variational value is $76.11\%$, at $\alpha = 0.90$ and threshold
$10^{-8}$;\ and at that same cell the $\sigma$-only solve returns $91.06\%$,
so opening the azimuthal channels there makes the answer \emph{worse} by
fifteen points.  The qualitative conclusion is therefore \emph{not} robust
pointwise across the grid, and this paper does not claim that it is.

What the result rests on instead, stated exactly:\ at the tight discard
thresholds, where the generalized eigenproblem is actually being solved rather
than regularized, every surviving variational point lies between $99.03$ and
$99.14\%$;\ and the independent Gaussian-basis route, which is well
conditioned and needs no discard threshold at all, gives $99.10\%$.  The claim
is the agreement of those two.  It is not a property of the whole
$(\alpha, \mathrm{threshold})$ grid, most of which is conditioning noise.  The $-1.161304$~Ha quoted in
Table~\ref{tab:convergence} sits $58\,\mu$Ha below the conditioned
value, i.e.\ the last two of its six decimals reflect linear
dependence rather than physics.  No conclusion in this paper turns
on them.  That the number is conditioning-\emph{limited} rather than
physics-limited is shown directly in Sec.~\ref{sec:recondition}:\
re-basing this same span removes the linear dependence and the energy
climbs to chemical accuracy.

\subsection{The monomial cap is conditioning, not a ceiling}
\label{sec:recondition}

\textbf{[MEASURED]} The $99.1\%$ above is the ceiling of the
\emph{monomial} radial factor $\xi^j$, not of the prolate geometry or
the electron-electron cusp.  The set $\{\xi^j e^{-\alpha\xi}\}$ is a
Hankel moment problem, exponentially ill-conditioned in the degree:\
the overlap reaches $\mathrm{cond} = 1.8\times10^{16}$ already at
$(3,3)$, $|m|\le1$ and $4.1\times10^{26}$ at $(5,5)$, so the basis
cannot be pushed further in double precision.  But the monomials
$\{\xi^j\}_{j\le J}$ and, e.g., the Laguerre polynomials
$\{L_j(2\alpha(\xi-1))\}_{j\le J}$ span the \emph{same} degree-$J$
space, so re-basing onto an orthogonal-polynomial family is an
\emph{exact} change of basis:\ it leaves the span---and hence the
exact-arithmetic energy---unchanged, and only the conditioning changes.

\textbf{[MEASURED]} Because the congruence $C S C^{\!\top}$ amplifies
the $10^{-16}$ entry error by $\|C\|^2$ and the high-degree monomial
matrices carry $\mathrm{cond}\sim10^{26}$, a float64 change of basis is
corrupted---it breaks for \emph{both} families at every truncation
tested.  Every matrix entry that passes through the monomial basis is
therefore built in extended precision,
the well-conditioned orthogonal matrices are downcast to float64, and
the generalized eigenproblem is solved on the unit-normalized
(correlation) matrices---a diagonal congruence that leaves the
eigenvalues unchanged but removes the norm-spread inflation of
$\mathrm{cond}$.  With that pipeline the energy climbs monotonically and
variationally past the monomial wall (Table~\ref{tab:recondition}):\
opening the $\delta$ channels on the re-based basis reaches $99.71\%$ of
$D_e$ at $(4,4)+\delta$ and $\gvq{p12_rebased_de_pct}{99.77}\%$
($\gvq{p12_rebased_err_mha}{0.41}$~mHa) at $(5,5)+\delta$---inside
chemical accuracy, a four-fold reduction of the residual from the
monomial $99.09\%$, every point containing its predecessor and lying
below it.

\textbf{[MEASURED]} That ladder is computed at a \emph{fixed} $\alpha = 1.0$,
which is what makes it a ladder:\ the statement that each point contains its
predecessor and lies below it holds only at a consistent exponent, so
Table~\ref{tab:recondition} keeps $\alpha$ fixed throughout.  But $\alpha$ is a
variational parameter, and $1.0$ is not its optimum.  Scanning it at the largest
truncation gives $0.406$, $0.346$, $\gvq{p12_rebased_err_mha_aopt}{0.32}$ and
$0.334$~mHa at $\alpha = 1.00$, $1.20$,
$\gvq{p12_rebased_alpha_opt}{1.40}$ and $1.50$, so the minimum is bracketed at
$\alpha = \gvq{p12_rebased_alpha_opt}{1.40}$ and the best variational result this
basis supports is $\gvq{p12_rebased_de_pct_aopt}{99.81}\%$ of $D_e$
($\gvq{p12_rebased_err_mha_aopt}{0.32}$~mHa)---better than the $\alpha = 1.0$
endpoint by $0.082$~mHa, a fifth of that residual.  The improvement is not a
trade:\ the optimal point is variational, retains all $1944$ functions, and is
\emph{more} tightly conditioned than $\alpha = 1.0$ ($4.9\times10^{10}$ against
$9.4\times10^{10}$).  Both numbers are therefore reported, and they are not
competing:\ $99.77\%$ is the endpoint of the fixed-$\alpha$ ladder, $99.81\%$ is
the basis's best variational point.

\begin{table}[t]
\caption{H$_2$ at $R = 1.4011$~bohr, re-based to orthogonal polynomials
  (Laguerre$\times$Legendre and the $\mu$-adapted Gegenbauer give the
  identical energy).  Same span as the monomial basis of
  Table~\ref{tab:azimuthal}, re-conditioned;\ $\alpha = 1.0$, the
  extended-precision $V_{ee}$ build and unit-normalized solve of the text.
  $D_e^{\rm exact} = 0.174475$~Ha.}
\label{tab:recondition}
\begin{ruledtabular}
\begin{tabular}{lrr}
  \multicolumn{1}{c}{basis / channels} & \multicolumn{1}{c}{$N$}
  & \multicolumn{1}{c}{$D_e\%$} \\
\hline
  monomial ceiling ($|m|\le1$) &  144 & 99.09 \\
  re-based $(3,3)$, $|m|\le1$  &  256 & 99.22 \\
  re-based $(4,4)+\delta$      &  975 & 99.71 \\
  re-based $(5,5)+\delta$      & 1944 & \gvq{p12_rebased_de_pct}{99.77} \\
\end{tabular}
\end{ruledtabular}
\end{table}

\textbf{[MEASURED]} The extended-precision requirement above constrains
the \emph{change of basis}, not the orthogonal matrices themselves, and so
does not apply to a matrix that is never built in the monomial basis at
all.  The one-body operators admit such a construction:\ every
one-electron block factorises as radial $\times$ angular, so $S$ and
$H_1$ can be assembled directly in the orthogonal basis from
one-dimensional blocks small enough to evaluate exactly, after which the
$O(N^2)$ two-electron assembly is ordinary double precision.  Built that
way they agree with the extended-precision route re-based to
$6\times10^{-16}$ of the matrix scale at $(5,5)+\delta$ in both families,
and every energy in Table~\ref{tab:recondition} is unchanged.  Extended
precision is therefore not a requirement of the one-body half at all:\ its
one-dimensional blocks are evaluated exactly because they are small enough
for that to cost nothing, not because a congruence forces it.  What does
still force it is the $V_{ee}$ re-basing, which accordingly accounts for
essentially the whole cost of the calculation.

\textbf{[MEASURED]} The residual $\sim\!0.4$~mHa is radial completeness
and the true electron-electron cusp, which converge as the slow
$L^{-3}$ partial-wave crawl---\emph{not} more azimuthal channels:\
$\varphi$ buys under $+0.06$~pp on top of $\delta$.  A literal $99.9\%$
therefore needs a much larger basis or explicit correlation
(geminals);\ a \emph{radial} re-basing cannot substitute for the
$e$--$e$ cusp.

\textbf{[MEASURED]} The $\mu$-weight-adapted family---the generalized
Laguerre $L_n^{(\mu)}$ for the $(\xi^2-1)^{\mu/2}$ weight and the
Gegenbauer $C_n^{(\mu+1/2)}$ for $(1-\eta^2)^{\mu/2}$, reducing to plain
Laguerre$\times$Legendre at $\mu = 0$---spans the identical space (the
energy agrees to $<1$~$\mu$Ha at every truncation;\ at $(5,5)+\delta$ and the
optimal $\alpha = \gvq{p12_rebased_alpha_opt}{1.40}$ the two families agree to
all seven printed digits, $E = -1.1741513$~Ha) at a normalized
condition number \emph{six} orders of magnitude smaller
($3.7\times10^4$ vs $4.9\times10^{10}$ at $(5,5)+\delta$, $\alpha = 1.40$;\
$9.1\times10^4$ vs $9.4\times10^{10}$ at $\alpha = 1.0$), which is what
keeps the downcast solve trustworthy at large truncation.
\textbf{[MEASURED]} Earlier versions of this sentence read ``two to three orders
of magnitude'' while quoting the $(5,5)+\delta$ pair, whose ratio is
$1.0\times10^{6}$;\ the ``two to three'' is the $326\times$ of the $(3,3)$,
$|m|\le1$ case and does not describe the figures it was attached to.  It does
\emph{not} reach a given accuracy with fewer functions:\ at equal degree
the span, and hence the energy, is the same---the gain is conditioning
alone.  A per-\emph{function} exponent instead (making the set
$L^2$-orthonormal) would break completeness and plateau the accuracy,
and is not used.

Backing:\ module \texttt{geovac/prolate\_recondition.py}, tests
\texttt{tests/test\_paper12\_recondition.py}.

%% ========================================================================
%% VII. IMPLICATIONS FOR THE GEOVAC PROGRAM
%% ========================================================================

\section{Implications for the GeoVac Program}
\label{sec:implications}

\subsection{The natural geometry hierarchy}

Papers~0, 1, and 7 established $S^3$ as the natural geometry
for one-center, one-electron systems.  Paper~11 established
prolate spheroids as the natural geometry for two-center,
one-electron systems.  This paper completes the analysis for
two-center, two-electron systems and identifies the hierarchy
of natural geometries:

\begin{enumerate}
  \item \textbf{One-center, one-electron} (H, He$^+$):
    Three-sphere $S^3$ via Fock projection.  The Schr\"odinger
    equation becomes a Laplace--Beltrami eigenvalue problem.
    Accuracy: $< 0.1\%$~\cite{loutey_paper0, loutey_paper1}.
  \item \textbf{Two-center, one-electron} (H$_2^+$):
    Prolate spheroidal coordinates $(\xi, \eta, \varphi)$.
    The Schr\"odinger equation separates into two coupled ODEs.
    Accuracy: 0.0002\% (spectral Laguerre)~\cite{loutey_paper11}.
  \item \textbf{Two-center, two-electron} (H$_2$):
    Prolate spheroidal $\times$ prolate spheroidal with
    Neumann-algebraic $V_{ee}$.  One-electron physics is
    exact; electron correlation is resolved to 92.4\% in the
    $\sigma$ sector and 99.1\% once the azimuthal channels are
    restored (Sec.~\ref{sec:azimuthal}).  \emph{This paper.}
\end{enumerate}

Each level in this hierarchy addresses a qualitatively different
physical singularity: the nuclear Coulomb singularity ($1/r$),
the two-center potential ($1/r_A + 1/r_B$), and the
electron-electron Coulomb singularity ($1/r_{12}$).  The GeoVac
program identifies each singularity's natural coordinate system and
discretizes the Laplace--Beltrami operator in that system.

\textbf{[WITHDRAWN]} A fourth entry appeared here in earlier
versions---``three-body coalescence: hyperspherical coordinates for
the interelectron cusp, identified but not yet implemented''---
presented as a consequence of this paper's residual.  It is removed,
because the residual was misdiagnosed (Sec.~\ref{sec:gap}).  This
is a statement about what \emph{this paper} licenses, not about
hyperspherical coordinates, whose use for two-electron atoms rests
on separate grounds in Paper~13~\cite{loutey_paper13}.  What the
present result removes is the \emph{evidence} that fourth entry
rested on, not a verdict in its favour.  Comparisons between the two
geometries are meaningful only at matched $m_{\max}$, and at matched
$m_{\max}$ neither this paper nor Paper~15~\cite{loutey_paper15}
claims an advantage over the other.  The two published figures are
not a matched pair:\ the $99.1\%$ here is
$(j_{\max},l_{\max}) = (3,3)$ at $|m| \le 1$, while the $96.0\%$
there is $l_{\max} = 6$ with a Schwartz cusp correction and a
different solver class, and the $99.97\%$ of
Ref.~\cite{tao_mccurdy_rescigno2010} is a grid method again
different from both.  The defensible statement is the negative one:\
prolate spheroidal coordinates are not disqualified for H$_2$ by an
inability to represent the cusp, because no such inability is in
evidence.

\subsection{Quantum chemistry without six-dimensional quadrature}

The Neumann algebraic $V_{ee}$ realizes a form of quantum
chemistry in which the two-electron integrals are reduced entirely to
recurrence relations and closed forms on quantum number labels---with
the log-singular second-kind moment $B_l$ now closed-form
(Sec.~\ref{sec:Bl_closed}), \emph{no} numerical integration remains in
either the $\sigma$ or the $|m| \ge 1$ sector.  Combined with the
analytical kinetic energy (integration by parts) and nuclear
attraction (simple polynomial moments), the \emph{entire}
Hamiltonian matrix is constructed without any quadrature at all.

This has conceptual significance beyond computational
efficiency: it demonstrates that the physics of
electron-electron repulsion in prolate spheroidal
coordinates is fully encoded in the algebraic structure
of the Legendre functions $P_l$ and $Q_l$, not in the
spatial geometry of a quadrature grid.  The Neumann expansion
is not an approximation; it is the exact representation of
$1/r_{12}$ in the prolate spheroidal basis, truncated only
by the angular momentum selection rule imposed by the basis
itself.

\subsection{Comparison with the \texorpdfstring{$S^3$}{S3}
approach}

The graph Laplacian FCI approach (Papers~0, 1, 7, and the
FCI papers) constructs $V_{ee}$ from Slater $F^0$ integrals
on the $S^3$ density overlap~\cite{loutey_paper7}, which are
computed analytically from hydrogen radial wavefunctions.
That approach also avoids numerical integration of $V_{ee}$
and achieves comparable accuracy for atoms.

The prolate spheroidal approach with Neumann $V_{ee}$ extends
this philosophy to molecules: the natural coordinate system
for each class of problems determines the algebraic structure
of the two-electron integrals.  On $S^3$, the algebraic
structure is the Slater expansion in radial quantum numbers.
On the prolate spheroid, it is the Neumann expansion in
Legendre orders.  Both are special cases of the general
principle: \emph{in the natural coordinate system, the
Coulomb kernel has a known multipole expansion that converts
integrals into algebraic sums.}

%% ========================================================================
%% VIII. CONCLUSION
%% ========================================================================

\section{Conclusion}
\label{sec:conclusion}

We have demonstrated that the Neumann expansion of $1/r_{12}$
in prolate spheroidal coordinates eliminates the electron-electron
integration bottleneck in the two-electron prolate spheroidal CI
for H$_2$.  The key results are:

\begin{enumerate}
  \item \textbf{92.4\% of exact $D_e$} with fully quadrature-free
    algebraic $V_{ee}$ (the log-singular $B_l$ moment closed-form,
    Sec.~\ref{sec:Bl_closed}) and zero fitted parameters.
  \item \textbf{12--20 percentage point improvement} over numerical
    $V_{ee}$ at every basis size from 1 to 72 functions.
  \item \textbf{Numerical $V_{ee}$ saturates at $\sim$80\%}
    regardless of basis size, confirming grid truncation error
    as the dominant error source.
  \item \textbf{27 basis functions is the sweet spot}: 92.2\%
    in 705~s versus 92.4\% in 5271~s at 72 functions.
  \item \textbf{The 7.6\% gap} is a basis-space limit, not an
    integration error---and specifically a \emph{configuration}
    limit:\ the $\varphi$-independent basis spans only
    $m_1 = m_2 = 0$, while a $^1\Sigma_g^+$ state constrains only
    $M = m_1 + m_2$.  The $\pi^2$ and $\delta^2$ configurations it
    excludes carry the angular correlation.
  \item \textbf{99.1\% of exact $D_e$} (stability envelope
    $99.0$--$99.1\%$) on restoring the azimuthal
    channels at $|m| \le 1$ in the same basis with the same
    Neumann kernel (Sec.~\ref{sec:azimuthal}), a gain of
    $\gvq{p12_azimuthal_gain_mha}{11.6}$~mHa, against
    $\gvq{p12_sigma_growth_mha}{0.34}$~mHa from near-tripling the $\sigma$
    basis.
    The $\sigma$-only value reproduces independently in a Gaussian
    basis to $0.2$~mHa, and the extended value to within the stated
    stability envelope.
  \item \textbf{$\gvq{p12_rebased_de_pct_aopt}{99.81}\%$ of exact $D_e$}
    ($\gvq{p12_rebased_err_mha_aopt}{0.32}$~mHa, chemical accuracy) once that
    $99.1\%$ monomial basis is re-conditioned and the shared exponent is taken to
    its variational optimum $\alpha = \gvq{p12_rebased_alpha_opt}{1.40}$
    ($\gvq{p12_rebased_de_pct}{99.77}\%$ at the fixed $\alpha = 1.0$ of
    Table~\ref{tab:recondition}):\ its $\xi^j$ radial factor
    is a Hankel moment problem (the $99.1\%$ is a conditioning wall, not
    a ceiling), and re-basing the same span onto orthogonal polynomials
    is an exact change of basis that lets the energy climb monotonically
    and variationally to chemical accuracy at $(5,5)+\delta$
    (Sec.~\ref{sec:recondition}).  The residual is radial completeness
    and the $e$--$e$ cusp, not more azimuthal channels.
\end{enumerate}

\textbf{[WITHDRAWN]} Earlier versions of this paper concluded that
the residual was the electron-electron cusp demanding non-analytic
$r_{12}^{1/2}$ and $r_{12}\ln r_{12}$ terms, and read that as
identifying hyperspherical coordinates as the next natural geometry.
Both are withdrawn:\ the calculations above contain no such term and
reach 99\%, and the grid-based prolate spheroidal treatment of
Ref.~\cite{tao_mccurdy_rescigno2010} reaches $0.05$~mHa in these
coordinates with a polynomial angular basis.

The result demonstrates a form of quantum chemistry where
recurrence relations and closed forms on quantum number labels replace
the six-dimensional quadrature for electron-electron repulsion
entirely---including the log-singular $B_l$ moment, whose closed form
(Sec.~\ref{sec:Bl_closed}) carries the same construction to the
associated Legendre functions of the $|m| \ge 1$ channels
(Sec.~\ref{sec:azimuthal}), so no numerical integration remains.
Just as
Paper~7 showed that atomic energies are encoded in the
topology of $S^3$, this paper shows that molecular
electron-electron repulsion is encoded in the algebraic
structure of the Legendre functions $P_l$, $Q_l$ and their
associated forms $P_l^m$, $Q_l^m$ on the
prolate spheroid.  The physics is in the algebra, not in the
grid.

%% ========================================================================
%% ACKNOWLEDGMENTS
%% ========================================================================

\begin{acknowledgments}
Computational resources were provided by personal hardware.
\end{acknowledgments}

%% ========================================================================
%% APPENDIX A: NUMERICAL VERIFICATION OF AUXILIARY INTEGRALS
%% ========================================================================

\appendix

\section{Numerical Verification of Auxiliary Integrals}
\label{app:verification}

All auxiliary integrals were validated against independent
numerical quadrature.  The test suite (\texttt{tests/test\_neumann\_vee.py},
32~tests) verifies:

\begin{itemize}
  \item \textbf{$C_l(q)$}: Base cases ($C_0$ even/odd, $C_1$ odd),
    selection rule, and known analytical values.  Agreement to
    machine precision ($< 10^{-14}$).
  \item \textbf{$A_n(\alpha)$}: Comparison with \texttt{scipy.integrate.quad}.
    Agreement to $< 10^{-12}$.
  \item \textbf{$A_l(p, \alpha)$}: Comparison with numerical
    quadrature of $\xi^p e^{-\alpha\xi} P_l(\xi)$.
    Agreement to $< 10^{-10}$.
  \item \textbf{$B_l(p, \alpha)$}: Comparison with independent
    adaptive quadrature.  Agreement to $< 10^{-10}$.
  \item \textbf{$X_l(p_1, p_2, \alpha)$}: Comparison with 2D
    numerical quadrature.  Agreement to $< 10^{-8}$.
    Symmetry $X_l(p_1, p_2) = X_l(p_2, p_1)$ verified to
    $< 10^{-14}$.
  \item \textbf{$V_{ee}$ matrix}: Symmetry, positive diagonal,
    and comparison with numerical elliptic integral quadrature.
    Neumann vs.\ numerical agreement to $< 0.01$~Ha per element
    (the discrepancy is the grid error in the numerical values).
  \item \textbf{Full Hamiltonian}: Variational principle (energy
    above exact), bound state, and consistency between
    Python and Numba code paths ($< 10^{-8}$~Ha).
\end{itemize}

\section{Comparison with \texorpdfstring{$r_{12}$}{r12}-Dependent
Basis Functions}
\label{app:r12}

We tested basis functions with explicit $r_{12}^p$ dependence
($p = 1, 2$), following the James--Coolidge
approach~\cite{James1933}.  The $r_{12}$ factor is expressed in
prolate spheroidal coordinates as a function of $(\xi_1, \eta_1,
\xi_2, \eta_2, \varphi_1 - \varphi_2)$, making the integrals
five-dimensional after azimuthal averaging.

Table~\ref{tab:r12_comparison} compares the best $r_{12}$-dependent
results with the $p = 0$ Neumann results.

\begin{table}[h]
\caption{Comparison of $r_{12}$-dependent basis (numerical
  $V_{ee}$) vs.\ $p = 0$ basis (Neumann $V_{ee}$).  All
  results at $R = 1.4011$~bohr.
  \label{tab:r12_comparison}}
\begin{ruledtabular}
\begin{tabular}{lddd}
  \multicolumn{1}{c}{Method}
  & \multicolumn{1}{c}{$N$}
  & \multicolumn{1}{c}{$D_e\%$}
  & \multicolumn{1}{c}{$t$ (s)} \\
\hline
  $r_{12}^1$, num., $60\times40$ grid & 6  & 85.7 & 447 \\
  $r_{12}^{0-2}$, num., $20\times14$ grid & 18 & 86.8 & 970 \\
  $p\!=\!0$, Neumann, $30\times20$ grid & 27 & 92.2 & 705 \\
  $p\!=\!0$, Neumann, $30\times20$ grid & 72 & 92.4 & 5271 \\
\end{tabular}
\end{ruledtabular}
\end{table}

The Neumann $p = 0$ result at 92.2\% exceeds the best
$r_{12}$-dependent numerical result (86.8\%) by 5.4~percentage
points, confirming that exact $V_{ee}$ evaluation is more
important than $r_{12}$-dependent basis functions when the
$V_{ee}$ integrals are computed numerically.  The combination
of $r_{12}$-dependent basis functions with Neumann-type
algebraic $V_{ee}$ is a natural extension left for future work.

%% ========================================================================
%% BIBLIOGRAPHY
%% ========================================================================

\begin{thebibliography}{30}

\bibitem{loutey_paper0}
J.~Loutey,
``Angular Momentum as Information Geometry: Isotropic Packing and
the Universal Angular Cross-Section,''
GeoVac Paper~0 (2026).

\bibitem{loutey_paper1}
J.~Loutey,
``The Geometric Atom: Atomic Structure as a Packing Problem,''
GeoVac Paper~1 (2026).

\bibitem{loutey_paper7}
J.~Loutey,
``The Dimensionless Vacuum: Recovering the Schr\"odinger Equation
from Scale-Invariant Graph Topology,''
GeoVac Paper~7 (2026).

\bibitem{loutey_paper11}
J.~Loutey,
``The Molecular Fock Projection: Diatomic Molecules as Prolate
Spheroidal Lattices,''
GeoVac Paper~11 (2026).

\bibitem{loutey_paper13}
J.~Loutey,
``The Hyperspherical Lattice: Two-Electron Atoms as Coupled
Channel Graphs,''
GeoVac Paper~13 (2026).

\bibitem{loutey_paper15}
J.~Loutey,
``The Level~4 Natural Geometry: Two-Center Two-Electron Molecules in
Molecule-Frame Hyperspherical Coordinates,''
GeoVac Paper~15 (2026).

\bibitem{tao_mccurdy_rescigno2010}
L.~Tao, C.~W. McCurdy, and T.~N. Rescigno,
``Grid-based methods for diatomic quantum scattering problems.
III. Double photoionization of molecular hydrogen in prolate
spheroidal coordinates,''
\emph{Phys. Rev. A} \textbf{82}, 023423 (2010).

\bibitem{Fock1935}
V.~Fock,
``Zur Theorie des Wasserstoffatoms,''
Z. Phys. \textbf{98}, 145--154 (1935).

\bibitem{Neumann1878}
C.~Neumann,
\textit{\"Uber die Entwicklung einer Function mit imagin\"aren
Argument nach den Kugelfunctionen erster und zweiter Art}
(Teubner, Leipzig, 1878).

\bibitem{Roothaan1951}
C.~C.~J. Roothaan,
``A study of two-center integrals useful in calculations on
molecular structure,''
J. Chem. Phys. \textbf{19}, 1445--1458 (1951).

\bibitem{Ruedenberg1951}
K.~Ruedenberg,
``A study of two-center integrals useful in calculations on
molecular structure.\ II.\ The two-center exchange integrals,''
J. Chem. Phys. \textbf{19}, 1459--1477 (1951).

\bibitem{Morse1953}
P.~M. Morse and H.~Feshbach,
\textit{Methods of Theoretical Physics}
(McGraw-Hill, New York, 1953).

\bibitem{James1933}
H.~M. James and A.~S. Coolidge,
``The ground state of the hydrogen molecule,''
J. Chem. Phys. \textbf{1}, 825--835 (1933).

\bibitem{Kolos1968}
W.~Ko\l{}os and L.~Wolniewicz,
``Improved theoretical ground-state energy of the hydrogen
molecule,''
J. Chem. Phys. \textbf{49}, 404--410 (1968).

\bibitem{Kato1957}
T.~Kato,
``On the eigenfunctions of many-particle systems in quantum
mechanics,''
Commun. Pure Appl. Math. \textbf{10}, 151--177 (1957).

\bibitem{Fock1954}
V.~Fock,
``On the Schr\"odinger equation of the helium atom,''
Izv. Akad. Nauk SSSR Ser. Fiz. \textbf{18}, 161--172 (1954).

\bibitem{Morgan1986}
J.~D. Morgan~III,
``Convergence properties of Fock's expansion for S-state
eigenfunctions of the helium atom,''
Theor. Chim. Acta \textbf{69}, 181--223 (1986).

\bibitem{Numba2015}
S.~K. Lam, A.~Pitrou, and S.~Seibert,
``Numba: A LLVM-based Python JIT compiler,''
in \textit{Proceedings of the Second Workshop on the LLVM
Compiler Infrastructure in HPC} (ACM, New York, 2015),
pp.\ 1--6.

\bibitem{QUADPACK}
R.~Piessens, E.~de Doncker-Kapenga, C.~W. \"Uberhuber,
and D.~K. Kahaner,
\textit{QUADPACK: A Subroutine Package for Automatic Integration}
(Springer-Verlag, Berlin, 1983).

\end{thebibliography}

\end{document}
